\documentclass[11pt]{article}
\usepackage[letterpaper,left=1in,right=1in,top=0.8in,bottom=0.8in]{geometry}
\usepackage[T1]{fontenc}
\usepackage{lmodern}
\usepackage[final]{microtype}
\usepackage[hidelinks]{hyperref}
\hypersetup{pdftitle={UCLA Mission Statement},pdfauthor={Wenrui Huang}}
\pagestyle{empty}
\setlength{\parindent}{0pt}
\setlength{\parskip}{0.65em}
\emergencystretch=1em
\begin{document}
\begin{center}
{\Large UCLA Mission Statement}\\[6pt]
{\large Wenrui Huang}
\end{center}
\vspace{2pt}

I would contribute to UCLA's mission through fundamental mathematical research,
careful teaching, and support for students developing confidence and independence
in mathematics. My education at the University of Science and Technology of China
(USTC) and my doctoral work at Brown have given me experience in different academic
settings. Teaching in both places has made me attentive to the steps an explanation
can leave unstated and to the responsibility of making advanced ideas accessible.

I have served as a teaching assistant for calculus at Brown and for topology and
geometry at USTC. In one topology conversation, a student was treating
a quotient space as a subset of the original space. Drawing two circles and pairing
their points helped us distinguish an equivalence class from a representative.
At Brown, I similarly use pictures of integration regions to help students explain
the bounds of an integral before calculating it. In calculus for physics and
engineering students, I also ask what an answer represents and whether its sign
and units make sense. These experiences have reinforced my view that teaching
should help students explain and evaluate their own reasoning.

At UCLA, I would build on these practices to support students with different
levels of preparation. During office hours, I ask students to identify the last
step they understand and work forward from there; before they leave, I ask them
to explain the approach in their own words. I would make expectations explicit,
provide specific feedback, and use recurring difficulties to improve my course
materials. Because mathematics opens pathways into many fields, helping students
progress through foundational courses is one concrete way to support UCLA's
commitment to educational opportunity and social mobility. I would also welcome
opportunities to guide students toward advanced courses and introductory reading
projects in analysis and differential equations.

My research addresses the long-time behavior of kinetic and dispersive partial
differential equations, including Vlasov and Hartree models, modified scattering,
and the stability of small BGK waves. These problems contribute to the mathematical
understanding of models of particle and wave dynamics. Through collaborative
projects, I have explored connections between kinetic and dispersive equations
and developed methods that apply across related models. At UCLA, I would pursue
this research, participate actively in seminars, and continue exchanging ideas
with collaborators at other institutions. I would share preprints through my
website and explain the questions motivating my work to audiences beyond my
immediate specialty, contributing to the dissemination of knowledge and UCLA's
international research community.

UCLA's connection with Los Angeles would also give me an opportunity to extend
these commitments beyond university courses. I would seek opportunities to
contribute to existing mathematics outreach activities, for example by preparing
a session on how pictures and simple examples help us reason about unfamiliar
objects. I would begin by learning from organizers about the needs of local
students and how my experience could be useful. My own experience moving between
academic settings would inform my efforts to welcome questions and make
participation in mathematical discussions easier.

Finally, I would support an effective department through dependable collaboration
with fellow instructors and teaching assistants, clear grading criteria, and
shared discussion of student difficulties. I would seek feedback on my teaching
and revise explanations that are not working. I hope to contribute to a community
in which research advances through open exchange and students receive the
guidance they need to pursue their ambitions.
\end{document}
